\section{Problems}\label{sec:problems}

Let V = \{1,\ldots{},n\} be a set of elements with symmetric distances $d_{ij}$ $\geq$ 0. A solution is a subset S $\subseteq$ V with |S| $\geq$ 2. We maximise two objectives,

\begin{equation*}
f_1(S) = \sum_{\substack{i<j\\ i,j\in S}} d_{ij}\;\; (\text{MaxSum}), \qquad f_2(S) = \min_{\substack{i<j\\ i,j\in S}} d_{ij}\;\; (\text{MaxMin}).
\end{equation*}

\noindent \textbf{MDP.} |S| = m, no further constraint. \textbf{CDP.} Each element has a capacity $c_{i}$ $\geq$ 0 and S must satisfy $\sum_{i\in S}$ $c_{i}$ $\geq$ B; |S| is free. In the benchmark B = $\mu$$\cdot$$\sum_{i\in V}$ $c_{i}$ with $\mu$ $\in$ \{0.2, 0.3\} (the suffixes b02 and b03 in the instance names). \textbf{GDP.} In addition each element has an opening cost $k_{i}$ and $\sum_{i\in S}$ $k_{i}$ $\leq$ K$_{1}$ (constant-cost model; the instance files also carry a second budget K$_{2}$ and unit costs of the variable-cost model, which we do not use, in line with the benchmark's own reference implementation). The instance files carry K$_{1}$ = $\nu$$\cdot$$\sum_{i\in V}$ $k_{i}$ with $\nu$ $\in$ \{0.2, 0.3\} (suffixes k02 and k03); Lozano-Osorio et al. tightened the budgets of the earlier benchmark of \citet{martinezgavara2021}, whose cost constraint was inactive at the optimum, so that the budget is active here. That the constant-cost reading is the one used for the published solutions is confirmed empirically: all 80 published n = 500 solutions are feasible under it (Section 5).

The original CDP and GDP literature optimises f$_{2}$ only, and in the GDP the number of selected elements is not fixed beforehand, which \citet{martinezgavara2021} consider more realistic than the cardinality of the classical MDP. Adding f$_{1}$ is our extension; it makes the feasible set's trade-off explicit: small sets raise MaxMin, large sets raise MaxSum. \textbf{The size of the set is a hidden dimension of this bi-objective problem.} Since f$_{1}$ never decreases when an element is added, any element at distance at least f$_{2}$(S) from every member of S can be inserted without lowering f$_{2}$ and with a larger f$_{1}$, so a solution is dominated by a larger one whenever such an element exists (and, on the GDP, fits in the budget). All methods in this study construct solutions up to the required capacity, exchange elements one by one and repair only what the constraints demand, so they work in the size regime induced by the capacity requirement; the fronts, the reference fronts and $\mathrm{HV}_{\mathrm{norm}}$ refer to that regime, and Section 6.13 quantifies what is left unexplored. The MDP has a fixed cardinality and is not affected.
